\section{Travaux Connexes}

\noindent\textcolor{titleblue}{\Large\bfseries L}'application du NLP au domaine juridique
a suscité un intérêt croissant ces dernières années. Cette section examine les
travaux existants dans trois axes : les systèmes de Q\&R juridique et RAG, les
modèles de langage et d'embedding pour l'arabe, et les stratégies de recherche
hybride.

\subsection{Systèmes RAG et Q\&R Juridique}

L'architecture RAG, introduite par \cite{lewis2020retrieval}, combine un
retriever (généralement Dense Passage Retrieval) avec un générateur
séquentiel (BART ou T5) pour produire des réponses contextualisées.
Depuis, plusieurs systèmes RAG spécialisés ont émergé dans le domaine
juridique. \cite{louis2023legal} présente Legal-RAG, un système de Q\&R
juridique pour le droit indien utilisant des embeddings multilingues et un
LLM fine-tuné. Leur approche atteint une précision de 72\% sur le jeu de
données Indian Legal Queries. \cite{nguyen2023vietnamese} développe un
système similaire pour le droit vietnamien, démontrant l'importance de
l'adaptation linguistique et culturelle dans les systèmes RAG juridiques.

\vspace{0.2cm}

\noindent Plus proche de notre contexte, \cite{alshareef2023arabic} propose
un système de Q\&R pour le droit saoudien utilisant AraBERT et Elasticsearch
pour la recherche lexicale. Leur système atteint une exactitude de 68\% sur
un jeu de 100 questions juridiques arabes. \cite{elaraby2024legal} présente
un système RAG pour le droit égyptien avec des embeddings multilingues et
un LLM distant. \cite{sharma2024indian} propose Indian Legal RAG, un
système multilingue pour le droit indien intégrant une recherche hybride.
\cite{huang2024sail} présente SAIL-Legal, un framework RAG scalable pour
le droit chinois avec indexation vectorielle distribuée.
\cite{zhou2023lexgpt} développe LexGPT, un assistant juridique basé sur
GPT-4 avec augmentation par recherche. Notre travail étend ces approches
en proposant une architecture complète, une recherche hybride avancée, et
un benchmark systématique sur le droit marocain.

\subsection{NLP Arabe pour le Domaine Juridique}

Le traitement automatique de l'arabe présente des défis bien documentés
\cite{habash2010arabic}. La morphologie riche et dérivationnelle de l'arabe
(racine trilitère, schèmes, affixes), la diglossie (MSA vs dialectes), et
l'orthographe ambiguë (alef sans hamza, tatweel, ligatures) rendent le
prétraitement linguistique crucial.

\vspace{0.2cm}

\noindent Dans le domaine juridique spécifiquement, \cite{zaghouani2021arabic}
souligne la complexité du vocabulaire juridique arabe, caractérisé par des
formules figées, des emprunts au français et à l'anglais, et une syntaxe
particulièrement élaborée. \cite{benali2025moroccan} note que les systèmes
de NLP arabe juridique marocain doivent gérer des phrases très longues (souvent plus
de 100 mots par article), un haut degré de co-référence, des
abréviations non standardisées, et un mélange fréquent d'arabe et de français
dans un même document. \cite{nassif2024arabic} propose une revue complète
des avancées récentes en NLP arabe, confirmant que la morphologie riche
de l'arabe et la diglossie restent des défis ouverts.

\vspace{0.2cm}

\noindent Un défi supplémentaire propre au contexte marocain est la
\textbf{qualité variable des documents sources}. La majorité des textes
de loi historiques existe uniquement au format PDF scanné (Bulletins
Officiels), avec des défauts de numérisation importants : inclinaison,
bruit de fond, caractères dégradés, et mélange de polices. Les techniques
de correction post-OCR par modèles de vision (VLM) \cite{benajiba2023arabic}
ont montré des résultats prometteurs pour l'arabe, avec une réduction des
erreurs de reconnaissance de 30 à 50\% sur des corpus juridiques. Notre
système adopte cette approche en cascade (EasyOCR $\rightarrow$ Gemma 4 VLM).

\vspace{0.2cm}

\noindent Des outils comme CAMeL Tools \cite{obaraj2021cameltools} offrent
des analyseurs morphologiques et des tagger POS pour l'arabe standard.
Cependant, leur application au texte juridique nécessite une adaptation,
le vocabulaire et les structures grammaticales juridiques différant
significativement de l'arabe journalistique ou littéraire. Dans notre système, nous avons opté pour une
approche de normalisation basée sur des expressions régulières et des
pondérations Unicode plutôt que des analyseurs morphologiques complets,
afin d'éviter la propagation d'erreurs et de maintenir un débit de
traitement élevé.

\subsection{Modèles d'Embedding Multilingues}

Le choix du modèle d'embedding est crucial dans tout système RAG. Pour
l'arabe, plusieurs études ont évalué les modèles multilingues.
\cite{ali2022arabic} compare LaBSE, XLM-R, et AraBERT sur des tâches de
similarité sémantique arabe, concluant que LaBSE offre les meilleures
performances cross-lingues. \cite{reimers2020multilingual} montre que les
modèles Sentence-BERT multilingues (mpnet-base, distiluse-base) transfèrent
bien vers l'arabe malgré un entraînement principalement sur des langues
européennes.

\vspace{0.2cm}

\noindent \cite{wang2023bge} introduit BGE-M3, un modèle d'embedding
multi-granularité (token, phrase, document) supportant plus de 100 langues,
dont l'arabe. \cite{liang2023multilingual} présente E5-large, un modèle
basé sur les paires texte-texte avec un pré-entraînement contrastif étendu.
Ces deux modèles (BGE-M3 et E5-large) produisent des embeddings de 1024
dimensions, offrant une capacité de représentation supérieure aux modèles
768 dimensions comme mpnet-base et LaBSE.

\vspace{0.2cm}

\noindent Notre benchmark sur le corpus juridique marocain (section 5) est
le premier, à notre connaissance, à comparer systématiquement ces cinq
modèles (E5-large, BGE-M3, LaBSE, mpnet-base, E5-base) avec des variantes
de requêtes multilingues (arabe formel, Darija, français, anglais) dans un
contexte juridique.

\subsection{Stratégies de Recherche Hybride}

La recherche dense (vectorielle) capture la similarité sémantique mais peut
manquer des correspondances lexicales exactes (numéros d'articles, termes
techniques spécifiques). La recherche BM25 excelle dans la correspondance
exacte mais ignore le contexte sémantique \cite{robertson2009bm25}. La
fusion de ces deux approches a montré des gains significatifs dans divers
domaines \cite{lin2021hybrid}.

\vspace{0.2cm}

\noindent Reciprocal Rank Fusion (RRF) \cite{cormack2009reciprocal} est une
méthode de fusion non paramétrique robuste qui combine les rangs de
plusieurs listes de classement. Contrairement à la fusion pondérée linéaire
(hybrid), RRF ne nécessite pas de calibration d'hyperparamètres sur un jeu
de validation et est moins sensible aux différences d'échelle entre les
scores des différents systèmes. Dans notre contexte juridique arabe, RRF
s'est avéré particulièrement adapté car les scores de similarité dense
(cosinus) et de BM25 (tf-idf normalisé) opèrent sur des échelles très
différentes.

\vspace{0.2cm}

\noindent \cite{shao2023hybrid} a démontré que la recherche hybride
Dense + BM25 + RRF offre les meilleurs résultats pour les requêtes
juridiques, confirmant les résultats de \cite{thakur2021hybrid} sur les
benchmarks BEIR et MTEB. Notre benchmark (section 5) confirme ces
résultats dans le domaine juridique arabe, avec une amélioration notable
du Hit@5 (68,75\% pour RRF contre 65,62\% pour dense et 46,88\% pour BM25).

\subsection{Large Language Models pour l'Arabe}

Les LLMs ouverts ont connu une adoption massive récemment. Gemma 4
\cite{google2025gemma4} est la quatrième génération des modèles ouverts
de Google, basée sur l'architecture Gemma. Disponible en tailles 2B et 9B,
elle supporte le multilingue dont l'arabe avec des performances compétitives.
Pour notre système, nous utilisons Gemma 4 2B en quantification Q4\_K\_XL via
llama.cpp \cite{gerganov2023llamacpp}, servie sur un GPU T4 via Google Colab.

\vspace{0.2cm}

\noindent Des études récentes \cite{hendy2023llm} montrent que les LLMs
multilingues comme Gemma et Llama peuvent répondre correctement à des
questions juridiques en arabe sans fine-tuning spécifique, à condition de
recevoir un contexte pertinent via RAG. L'approche \textit{zero-shot} avec
injection de contexte (grounding) est préférée au fine-tuning car elle évite
le surapprentissage et permet de maintenir la base de connaissances
juridiques à jour sans réentraînement.

\subsection{Synthèse Comparative}

Le tableau~\ref{tab:related_synthesis} synthétise les caractéristiques
techniques des principaux systèmes RAG juridiques récents comparés à
notre proposition.

\begin{table}[H]
\centering
\caption{Synthèse comparative des systèmes RAG juridiques récents}
\label{tab:related_synthesis}
\renewcommand{\arraystretch}{1.15}
\resizebox{\columnwidth}{!}{%
\begin{tabular}{lcccc}
\toprule
\textbf{Système} & \textbf{Dom.} & \textbf{Rech.} & \textbf{Embedding} & \textbf{P. max} \\
\midrule
Lewis+ (2020) & Général & DPR & BERT & 68\%\\
\cite{lewis2020retrieval} & & & & (R@20)\\
\midrule
Louis+ (2023) & Indien & Dense & e5 & 72\%\\
\cite{louis2023legal} & & (FAISS) & & (P@1)\\
\midrule
Alshareef+ (2023) & Saoudien & Lexical & AraBERT & 68\%\\
\cite{alshareef2023arabic} & & (ES) & & (Acc.)\\
\midrule
Zhou+ (2023) & US & Hybride & Ada-002 & N/D\\
\cite{zhou2023lexgpt} & & (GPT-4) & &\\
\midrule
Nguyen+ (2023) & Vietnam. & Dense & PhoBERT & 65\%\\
\cite{nguyen2023vietnamese} & & (FAISS) & & (Acc.)\\
\midrule
Elaraby+ (2024) & Égyptien & Hybride & E5-l. & 70\%\\
\cite{elaraby2024legal} & & (RRF) & & (F1)\\
\midrule
Sharma+ (2024) & Indien & Hybride & MuRIL & 73,5\%\\
\cite{sharma2024indian} & & (RRF) & & (MRR)\\
\midrule
Huang+ (2024) & Chinois & Hybride & BGE & 76,2\%\\
\cite{huang2024sail} & & (Dist.) & & (H@5)\\
\midrule
\textbf{Notre sys.} & \textbf{Maroc} & \textbf{Hybr.} & \textbf{E5-l.} & \textbf{81\%}\\
\textbf{(2026)} & \textbf{7 dom.} & \textbf{RRF} & \textbf{1024d} & \textbf{MRR}\\
\bottomrule
\end{tabular}
}%
\end{table}

\vspace{0.2cm}

\noindent Notre système se distingue par plusieurs aspects : (1) la
combinaison de la recherche hybride RRF avec un tokenizer arabe spécialisé
pour le canal BM25 ; (2) l'évaluation systématique de 5 modèles d'embedding
multilingues sur un corpus juridique arabe avec 65 variantes de requêtes ;
(3) un pipeline complet de la numérisation OCR à la génération de réponses,
totalement open-source et conteneurisé. Seul le système de \cite{elaraby2024legal}
partage une architecture similaire pour le monde arabe, mais notre corpus
est plus étendu (2\,357 articles vs. 800) et notre benchmark plus approfondi.